\vspace{-1mm}
\section{Related Work}
\label{sec:related}
\noindent\textbf{Theory of activation steering.}~Several recent works have sought to develop a unified theory of activation steering by interpreting interventions as structured transformations in representation space. Linear Activation Transportation~\citep{rodriguez2025controlling} formulates steering as a transport map between source and target activation distributions, LiSeCo~\citep{cheng2024linearly} analyzes steering through linear separability and constrained optimization, and ODESteer~\citep{zhaoodesteer} views steering as a continuous-time dynamical process rather than a one-shot activation edit. \citet{taimeskhanov2026towards} theoretically study the effect of steering strength in additive interventions, characterizing its impact on next-token probabilities, concept expression, and cross-entropy. These perspectives provide useful formal tools for reasoning about activation interventions. These perspectives provide useful formal tools for reasoning about activation interventions. However, their guarantees are typically derived for surrogate objectives and under specific geometric constraints rather than from a principled perspective. Also, the Schr\"odinger Bridge (SB) problem has been widely adopted for generative modeling via forward--backward SDE solvers~\citep{de2021diffusion, chen2021likelihood, shi2023diffusion}. \Cobras is the first to apply this
framework for LLM activation steering.

\noindent\textbf{Activation steering approaches.}~Early and widely used steering approaches~\cite{zou2023representation, li2023inference, rimsky2024steering} identify behavior-relevant directions in activation space and apply additive edits. More recent methods refine this basic recipe by accounting for activation geometry or distributional structure, for example through minimally invasive affine maps~\citep{singh2024representation}, norm-preserving rotations~\citep{pham2024householder}, transport-based maps~\citep{rodriguez2025controlling}, hyperspherical interventions~\citep{you2026spherical}, or continuous dynamical updates~\citep{zhaoodesteer}. Despite these advances, most existing methods are query-independent: once a direction or transformation is learned, essentially the same intervention is applied across prompts. Prior work shows that such fixed interventions can have limited OOD generalization, induce off-target side effects, and degrade performance when steering is unnecessary~\citep{tan2024analysing,raedler2025necessity,da2025steering}. Our work instead proposes a query-adaptive steering method that modulates the intervention based on the input, preserving OOD performance while improving in-distribution performance.
 

% Representation Engineering (RepE)~\citep{zou2023representation} showed that high-level concepts and behaviors can often be represented as directions in hidden-state/activation space, enabling direct manipulation of model representations. Building on this idea, Inference-Time Intervention (ITI)~\citep{li2023inference} steers selected attention-head activations along truthfulness-related directions, while Contrastive Activation Addition (CAA)~\citep{rimsky2024steering} constructs steering vectors from differences between positive and negative examples and adds them to residual-stream activations during generation. More recent methods move beyond simple additive interventions by considering richer activation geometry and distributional structure: Minimally Modified Counterfactuals (MiMiC)~\citep{singh2024representation} applies minimally invasive affine transformations to match counterfactual activation statistics, Householder Pseudo-Rotation (HPR)~\citep{pham2024householder} uses norm-preserving pseudo-rotations, Linear Activation Transport (Linear-AcT)~\citep{rodriguez2025controlling} formulates steering as a transport map between source and target activation distributions, SphericalSteering~\citep{you2026spherical} performs steering on the hypersphere to preserve activation magnitudes, and ODESteer~\citep{zhaoodesteer} interprets steering as a continuous dynamical process rather than a one-shot edit.

% Most of the present steering methods, are query indepdent, they end up steering using a constant steering direction which leads to catastrophic failure of performance in OOD tasks. In our work, we propose a query adaptive steering method that retains it performacne in OOD tasks while improving perofrmance on in-diostribution tasks. 

% Query indepdent ; Cobras is query adaptive ; see some failure cases; \textcolor{blue}{AD: mention catastrophic performance of OOD as the result of not being query-adaptive} 
